\section{Ciele riešenia}
\label{sec:ciele}

\subsection{Hlavný cieľ}

Hlavným cieľom projektu je vytvoriť prototyp adaptívneho systému, ktorý
z~dostupných podkladov o~AI modeli alebo AI systéme automatizovane vytvorí
jeho bezpečnostný a~governance profil, identifikuje riziká, nedostatky
a~chýbajúce informácie a~každé zistenie prepojí s~dôkazom, z~ktorého vzniklo.

\subsection{Čiastkové ciele}

Čiastkové ciele zodpovedajú piatim úlohám oficiálneho zadania
(kapitola~\ref{sec:zadanie}):

\begin{enumerate}
    \item \textbf{Jednotný informačný model.} Navrhnúť štruktúru, ktorá
    jednotne opisuje AI riešenie v~oblastiach model card, účel a~obmedzenia
    použitia, informácie o~dátach, evaluačné výsledky, bezpečnostné
    vlastnosti a~ďalšie relevantné charakteristiky.
    \item \textbf{Zber a~extrakcia informácií.} Implementovať automatizované
    získavanie informácií z~model cards, technickej dokumentácie a~ďalších
    zdrojov, ich extrakciu do informačného modelu a~vzájomné prepojenie.
    \item \textbf{Hodnotiace mechanizmy.} Implementovať vybrané automatizované
    hodnotenia: tvorbu alebo dopĺňanie model card, threat modeling,
    identifikáciu bezpečnostných rizík a~kontrolu splnenia definovaných
    požiadaviek.
    \item \textbf{Adaptívne hodnotenie.} Umožniť systému rozpoznať chýbajúce
    informácie, generovať doplňujúce otázky a~ku každému zisteniu uvádzať
    mieru istoty.
    \item \textbf{Vyhodnotenie systému.} Overiť systém na vybraných AI modeloch
    alebo systémoch z~hľadiska správnosti a~úplnosti získaných informácií,
    kvality hodnotení, identifikovaných rizík a~dohľadateľnosti zistení
    k~zdrojovým dôkazom.
\end{enumerate}
